\chapter*{Notación}
\addcontentsline{toc}{chapter}{Notación}

\section*{Conjuntos y estructuras}

\begin{tabularx}{\linewidth}{l X}
  \toprule
  Símbolo & Significado \\
  \midrule
  $\NN$, $\ZZ$, $\RR$ & naturales, enteros, reales \\
  $\BB$ & $\{\bot, \top\}$, valores booleanos \\
  $\Set$ & categoría de conjuntos y funciones \\
  $\Graph$ & categoría de grafos dirigidos y morfismos \\
  $\TGraph$ & categoría de grafos \emph{tipados} (la que vive sotFS) \\
  $G \xrightarrow{f} H$ & morfismo de grafos $f$ entre $G$ y $H$ \\
  $A \push_C B$ & pushout de $A$ y $B$ a lo largo de $C$ \\
  $|S|$ & cardinalidad del conjunto $S$ \\
  $\mathcal{P}(S)$ & conjunto potencia de $S$ \\
  \bottomrule
\end{tabularx}

\section*{TypeGraph (capítulo \ref{cap:typegraph})}

\begin{tabularx}{\linewidth}{l X}
  \toprule
  Símbolo & Significado \\
  \midrule
  $\NInode$ & nodo de tipo inodo (metadata POSIX) \\
  $\NDir$ & nodo de tipo directorio \\
  $\NCap$ & nodo de tipo capability \\
  $\NTxn$ & nodo de tipo transaction \\
  $\NVer$ & nodo de tipo version (snapshot) \\
  $\NBlk$ & nodo de tipo block (datos) \\
  $\econtains, \egrants, \edelegates$ & aristas tipadas \\
  $\ederivedfrom, \esupersedes, \epointsto, \ehasxattr$ & aristas tipadas \\
  $G$ & un \emph{TypeGraph} concreto (instancia) \\
  $\hashof{x}$ & hash de contenido \blake{} de $x$ \\
  \bottomrule
\end{tabularx}

\section*{DPO}

\begin{tabularx}{\linewidth}{l X}
  \toprule
  Símbolo & Significado \\
  \midrule
  $\dpo{L}{K}{R}$ & producción DPO con \emph{left}, \emph{glue}, \emph{right} \\
  $m: L \to G$ & \emph{matching} de la regla en el host $G$ \\
  $G \xRightarrow{\rho, m} H$ & paso DPO de $G$ a $H$ vía $(\rho,m)$ \\
  $\op{nac}$ & \emph{negative application condition} \\
  \bottomrule
\end{tabularx}

\section*{Operaciones POSIX y errno}

Las operaciones POSIX se denotan en \texttt{monospace}: \opcreate, \opmkdir,
\opunlink, \oprename, \oplink, \oprmdir, \opwrite, \opread. Los errores
POSIX se citan como \eexist, \enoent, \enotempty, \eacces, \eloop.

\section*{Verificación}

\begin{tabularx}{\linewidth}{l X}
  \toprule
  Símbolo & Significado \\
  \midrule
  TLA\textsuperscript{+} & lenguaje de specs de Lamport \cite{Lamport2002} \\
  TLC & el model checker de TLA\textsuperscript{+} \\
  Coq & asistente de pruebas \cite{Bertot2004} \\
  \texttt{Admitted} & lema enunciado pero no probado en Coq \\
  \citaTLA{x.tla} & cita al spec \texttt{x.tla} en \texttt{formal/} \\
  \citaCoq{x.v} & cita al archivo \texttt{x.v} en \texttt{formal/coq/} \\
  \citaCrate{x} & cita al crate \texttt{x} en el workspace \\
  \citaCommit{sha} & cita al commit \texttt{sha} del repo \\
  \bottomrule
\end{tabularx}

\section*{Cajas}

\begin{tabularx}{\linewidth}{l X}
  \toprule
  Tipo & Uso \\
  \midrule
  \emph{Definición} & introducción de un objeto formal \\
  \emph{Teorema} & resultado central, probado o citado \\
  \emph{Lema} & resultado auxiliar \\
  \emph{Invariante} & propiedad mantenida por construcción \\
  \emph{Trampa} & error histórico del proyecto + lección \\
  \emph{Observación} & nota lateral, no crítica para la lectura \\
  \emph{Ejercicio} & propuesta al lector, graduada $\star$/$\star\star$/$\star\star\star$ \\
  \bottomrule
\end{tabularx}

\clearpage
